\chapter*{Заключение}
\addcontentsline{toc}{chapter}{Заключение}

В ходе выполнения выпускной квалификационной работы была разработана и исследована программная система траекторного управления квадрокоптером в трёхмерном пространстве. Аналитический алгоритм согласованного управления по выходу был адаптирован под численное моделирование и дополнен интеллектуальным модулем выбора параметрической скорости движения вдоль траектории.

В работе удалось получить следующие основные результаты:
\begin{enumerate}
    \item Сформулирована постановка задачи траекторного управления нелинейным объектом с учётом геометрии пространственной кривой и введены ошибки слежения в системе координат Френе.
    \item Реализована нелинейная модель квадрокоптера с интеграторными расширениями и базовый аналитический регулятор траекторного слежения, проверенный на наборе тестовых кривых (прямая, окружность, спираль).
    \item Разработана процедура формирования датасета: для каждой точки кривой с помощью оракульного перебора определялось максимальное значение $V^{\ast}$, сохраняющее качественное слежение в коротком rollout.
    \item Обучено четыре модели выбора скорости с принципиально разными подходами к аппроксимации: базовая нейросетевая регрессия (MLP), стохастический актор с максимизацией энтропии (SAC), детерминированный актор с двойными критиками (TD3) и политика с ограниченным шагом обновления (PPO).
    \item Решён ряд технических проблем интеграции: устранён разрыв опорного сигнала при адаптивном изменении $V^{\ast}$, введён монитор безопасности против накопления $\eta$-интегратора, исправлена инициализация наблюдателя.
    \item Проведены сравнение и компьютерное моделирование моделей, демонстрирующие корректность работы всей системы в целом.
\end{enumerate}

Сравнительный анализ показал, что разные архитектуры по-разному балансируют между скоростью прохождения траектории и устойчивостью слежения.
MLP обеспечивает умеренный и стабильный прирост средней скорости (около $1.3\times$ относительно baseline) и хорошо работает на траекториях, близких к обучающему датасету.
SAC и TD3 позволяют достичь более высоких скоростных показателей при этом сохраняя устойчивость движения.
Наиболее сбалансированный результат на полном наборе бенчмарков показала модель SAC: ускорение без заметной деградации качества слежения.

PPO в части сценариев даёт наибольший прирост средней скорости, однако качество слежения оказывается менее стабильным.
TD3 занимает промежуточное положение между SAC и PPO и хорошо балансирует между ускорением и точностью движения.
Полученные результаты также подтвердили, что supervised MLP может быть вполне конкурентоспособен по отношению к более сложным RL-подходам.

Принципиальным результатом работы является двухуровневая схема управления: аналитический регулятор отвечает за устойчивость движения, а обучаемый модуль - только за выбор
скорости.
Это позволяет менять и улучшать интеллектуальную надстройку, без изменения структуры базового регулятора.

Дальнейшее развитие работы может быть связано с расширением набора тестовых траекторий, увеличением объёма датасета, исследованием других вариантов функции награды, а также
переносом системы в среду с реалистичной аэродинамикой и сенсорным шумом.